\documentclass[9pt,landscape,a4paper]{article}
\usepackage[utf8]{inputenc}
\usepackage[T1]{fontenc}
\usepackage{geometry}
\usepackage{multicol}
\usepackage{amsmath,amssymb}
\usepackage{enumitem}
\usepackage{graphicx}
\usepackage{xcolor}
\usepackage{listings}
\usepackage{booktabs}
\usepackage{titlesec}
\usepackage{fancyhdr}
\usepackage{hyperref}
\usepackage{CJKutf8}

% Page layout
\geometry{top=0.7cm, bottom=0.7cm, left=0.7cm, right=0.7cm}
\pagestyle{empty}
\setlength{\parindent}{0pt}
\setlength{\parskip}{1.5pt}
\setlength{\columnsep}{7pt}

% Section formatting
\titleformat{\section}{\normalsize\bfseries\color{blue!70!black}}{}{0em}{}[\titlerule]
\titleformat{\subsection}{\small\bfseries\color{blue!50!black}}{}{0em}{}
\titlespacing{\section}{0pt}{4pt}{2pt}
\titlespacing{\subsection}{0pt}{3pt}{1pt}

% Code listing style
\lstset{
    language=Python,
    basicstyle=\ttfamily\fontsize{6.5}{7.5}\selectfont,
    keywordstyle=\color{blue!70!black}\bfseries,
    stringstyle=\color{red!60!black},
    commentstyle=\color{green!50!black}\itshape,
    numberstyle=\tiny\color{gray},
    backgroundcolor=\color{gray!8},
    frame=single,
    framerule=0.3pt,
    rulecolor=\color{gray!50},
    breaklines=true,
    breakatwhitespace=false,
    tabsize=2,
    showstringspaces=false,
    aboveskip=2pt,
    belowskip=2pt,
    xleftmargin=2pt,
    xrightmargin=2pt,
}

\newcommand{\keybox}[1]{%
    \begin{center}
    \fcolorbox{blue!50!black}{blue!5}{%
        \parbox{0.95\linewidth}{\small #1}%
    }
    \end{center}
}

\newcommand{\zhint}[1]{{\scriptsize\color{gray!70!black}#1}}

\begin{document}
\begin{CJK}{UTF8}{gbsn}

\begin{center}
{\large\bfseries\color{blue!70!black} Chapter 7 -- Data Modeling II}\\
{\footnotesize IERG4320/IEMS5726 Data Science in Practice | Cheat Sheet}
\end{center}

\begin{multicols}{3}

% ==================== SECTION 1 ====================
\section{Clustering Overview}
\zhint{聚类：无监督学习，把相似的数据分到同一组，没有标签}

\begin{itemize}[leftmargin=8pt,itemsep=0pt,topsep=0pt]
\item Find groupings where data in each group is similar
\item Unsupervised -- no labels needed
\item Uses: basket analysis, news clustering, social network communities, dimensionality reduction
\end{itemize}

\textbf{Top 3 Algorithms}: K-Means, Agglomerative, Spectral

% ==================== SECTION 2 ====================
\section{K-Means Clustering}
\zhint{K-Means：指定K个簇，迭代移动中心点直到收敛}

\subsection{Algorithm}
\begin{enumerate}[leftmargin=10pt,itemsep=0pt,topsep=0pt]
\item Initialize k centroids (random guesses)
\item Assign each point to nearest centroid
\item Recalculate centroid = mean of cluster
\item Repeat 2--3 until no change
\end{enumerate}
\textbf{Parameters}: \texttt{k} (num clusters), \texttt{n\_init} (num initializations)

\subsection{Code: K-Means}
\begin{lstlisting}
from sklearn.cluster import KMeans
from sklearn.preprocessing import StandardScaler
from sklearn.datasets import make_blobs
from sklearn.metrics import silhouette_score
feat, true_labels = make_blobs(
  n_samples=200, centers=3,
  cluster_std=2.75, random_state=42)
scaler = StandardScaler()
scaled = scaler.fit_transform(feat)
kmeans = KMeans(init="random", n_clusters=3,
  n_init=10, max_iter=300, random_state=42)
kmeans.fit(scaled)
print("SSE:", kmeans.inertia_)
print("Centroids:", kmeans.cluster_centers_)
print("Labels:", kmeans.labels_[:5])
\end{lstlisting}
\zhint{inertia\_=SSE(越小越好); cluster\_centers\_=中心点坐标}

\subsection{Elbow Method -- Find Best k}
\zhint{肘部法则：SSE随k下降，拐点处就是最佳k}
\begin{lstlisting}
from kneed import KneeLocator
sse = []
for k in range(1, 11):
  km = KMeans(n_clusters=k, init="random",
    n_init=10, max_iter=300, random_state=42)
  km.fit(scaled)
  sse.append(km.inertia_)
# Plot SSE vs k, find elbow
kl = KneeLocator(range(1,11), sse,
  curve="convex", direction="decreasing")
print(kl.elbow)  # e.g. 3
\end{lstlisting}

% ==================== SECTION 3 ====================
\section{Agglomerative Clustering}
\zhint{层次聚类(自底向上)：每个点是一个簇，逐步合并最近的两个}

\subsection{Algorithm}
\begin{enumerate}[leftmargin=10pt,itemsep=0pt,topsep=0pt]
\item Start: each point = one cluster
\item Find two closest clusters, merge them
\item Update proximity matrix
\item Repeat until all merged into one cluster
\end{enumerate}

\textbf{Linkage types}: \zhint{计算簇间距离的方式}
\begin{itemize}[leftmargin=8pt,itemsep=0pt,topsep=0pt]
\item \textbf{Single}: min distance between clusters
\item \textbf{Complete}: max distance between clusters
\item \textbf{Average}: mean distance between clusters
\item \textbf{Ward}: minimize error sum of squares \zhint{最常用}
\end{itemize}

\subsection{Code: Agglomerative Clustering}
\begin{lstlisting}
from sklearn import cluster
from sklearn.neighbors import kneighbors_graph
# Ward linkage on Gaussian blobs
hc = cluster.AgglomerativeClustering(
  n_clusters=4, metric='euclidean',
  linkage='ward').fit_predict(dataset1)
# With connectivity constraints
connect = kneighbors_graph(
  dataset2, n_neighbors=5, include_self=False)
hc_conn = cluster.AgglomerativeClustering(
  n_clusters=2, metric='euclidean',
  linkage='complete',
  connectivity=connect).fit_predict(dataset2)
\end{lstlisting}
\zhint{connectivity=限制只能合并相邻的点，适合非球形簇}

% ==================== SECTION 4 ====================
\section{Spectral Clustering}
\zhint{谱聚类：用相似度矩阵的特征值降维后再聚类，适合复杂形状}

\subsection{Algorithm}
\begin{enumerate}[leftmargin=10pt,itemsep=0pt,topsep=0pt]
\item Build distance matrix $\to$ affinity matrix A
\item Compute degree matrix D, Laplacian L = D -- A
\item Find eigenvalues/eigenvectors of L
\item Use top k eigenvectors as new features
\item Apply K-Means in this new space
\end{enumerate}

\subsection{Code: Spectral Clustering}
\begin{lstlisting}
from sklearn.cluster import SpectralClustering
from sklearn.datasets import make_blobs
from numpy import random
random.seed(1)
x, _ = make_blobs(n_samples=400,
  centers=4, cluster_std=1.5)
sc = SpectralClustering(n_clusters=4).fit(x)
labels = sc.labels_
# Try different cluster numbers
for i in range(2, 6):
  sc = SpectralClustering(n_clusters=i).fit(x)
  plt.scatter(x[:,0],x[:,1],s=5,c=sc.labels_)
\end{lstlisting}
\zhint{默认affinity='rbf'; assign\_labels='kmeans'}

% ==================== SECTION 5 ====================
\section{Clustering Evaluation Metrics}
\zhint{聚类没有标签，用内部指标评估簇的质量}

\subsection{1. Silhouette Score}
\zhint{轮廓系数：衡量簇内紧密+簇间分离，范围[-1,1]，越大越好}\\
$s(i) = \frac{b(i) - a(i)}{\max(a(i), b(i))}$\\
\textbf{a(i)}: avg distance to points in same cluster \zhint{越小越好}\\
\textbf{b(i)}: avg distance to nearest other cluster \zhint{越大越好}

\begin{lstlisting}
from sklearn.metrics import silhouette_score
km = KMeans(n_clusters=3, init='k-means++',
  max_iter=300, n_init=10, random_state=0)
km.fit_predict(X)
score = silhouette_score(X, km.labels_,
  metric='euclidean')
print('Silhouette:', score)  # e.g. 0.553
# Visualize with yellowbrick
from yellowbrick.cluster import SilhouetteVisualizer
vis = SilhouetteVisualizer(km,colors='yellowbrick')
vis.fit(X)
\end{lstlisting}
\zhint{图中宽度=簇质量，高度=簇大小}

\subsection{2. Calinski-Harabasz Index}
\zhint{CH指数：簇间离散度/簇内离散度，越大越好}\\
$CH = \frac{B_k / (K-1)}{W_k / (N-K)}$\\
$B_k$: between-cluster dispersion; $W_k$: within-cluster.

\begin{lstlisting}
from sklearn import metrics
labels = kmeans.labels_
ch = metrics.calinski_harabasz_score(X, labels)
print(ch)  # e.g. 561.63, higher=better
\end{lstlisting}

\subsection{3. Davies-Bouldin Index}
\zhint{DB指数：簇内散度与簇间距离之比，越小越好}\\
$R_{i,j} = \frac{S_i + S_j}{M_{i,j}}$, \quad $DB = \frac{1}{K}\sum \max_{j \neq i} R_{i,j}$

\begin{lstlisting}
from sklearn.metrics import davies_bouldin_score
db = davies_bouldin_score(X, labels)
print(db)  # e.g. 0.768, lower=better
# Compare different k values
for i in range(2, 11):
  km = KMeans(n_clusters=i, random_state=30)
  labels = km.fit_predict(X)
  print(i, davies_bouldin_score(X, labels))
\end{lstlisting}

\keybox{
\textbf{Metric Summary}:\\
\textbf{Silhouette}: $[-1, 1]$, \textbf{higher} = better \zhint{轮廓系数}\\
\textbf{Calinski-Harabasz}: $[0, \infty)$, \textbf{higher} = better \zhint{CH指数}\\
\textbf{Davies-Bouldin}: $[0, \infty)$, \textbf{lower} = better \zhint{DB指数}
}

% ==================== SECTION 6 ====================
\section{Reinforcement Learning (RL)}
\zhint{强化学习：智能体在环境中通过试错获得奖励来学习策略}

\subsection{Key Concepts}
\begin{itemize}[leftmargin=8pt,itemsep=0pt,topsep=0pt]
\item \textbf{Agent}: the learner / decision maker \zhint{智能体}
\item \textbf{Environment}: what agent interacts with \zhint{环境}
\item \textbf{State} $s$: current situation \zhint{状态}
\item \textbf{Action} $a$: what agent can do \zhint{动作}
\item \textbf{Reward} $r$: feedback signal \zhint{奖励}
\item \textbf{Episode}: sequence of steps from start to end \zhint{一局游戏}
\item \textbf{Policy} $\pi(s) = a$: strategy for choosing actions \zhint{策略}
\end{itemize}

\textbf{Environment types}:\\
Discrete (grid, board game) vs Continuous (real world)\\
\zhint{离散环境=有限选择; 连续环境=无限可能}

\subsection{Exploration vs Exploitation}
\zhint{探索=尝试新动作; 利用=用已知最优动作}
\begin{itemize}[leftmargin=8pt,itemsep=0pt,topsep=0pt]
\item \textbf{Exploration}: try new actions (may find better) \zhint{探索}
\item \textbf{Exploitation}: use known best action \zhint{利用}
\item \textbf{$\varepsilon$-greedy}: with prob $\varepsilon$ explore randomly, else exploit
\item \textbf{Boltzmann}: explore more when getting negative rewards
\end{itemize}

\subsection{Q-Learning}
\zhint{Q-Learning：用Q表存储每个(状态,动作)的价值}\\
\textbf{Q-Table}: rows = states, cols = actions, values = expected future rewards.\\[2pt]
\textbf{Update formula}:\\
$Q(S,A) \leftarrow Q(S,A) + \alpha\big[R + \gamma \max_{A'} Q(S',A') - Q(S,A)\big]$\\[2pt]
$\alpha$ = learning rate, $\gamma$ = discount factor, $R$ = reward.\\
\zhint{当前Q值 + 学习率 $\times$ (实际回报 - 当前估计)}

\subsection{Code: Q-Learning Agent (Tic-Tac-Toe)}
\begin{lstlisting}
class QLearningAgent:
  def __init__(self, alpha=0.1, gamma=0.9,
               epsilon=0.1):
    self.q_table = defaultdict(
      lambda: np.zeros(9))
    self.alpha = alpha    # learning rate
    self.gamma = gamma    # discount factor
    self.epsilon = epsilon # exploration rate

  def get_action(self, state, available):
    state_key = self._state_to_key(state)
    if random.random() < self.epsilon:
      return random.choice(available) # Explore
    else:
      q = self.q_table[state_key]     # Exploit
      avail_q = {m: q[m] for m in available}
      return max(avail_q, key=avail_q.get)

  def update_q_value(self, state, action,
                     reward, next_state,
                     next_avail):
    sk = self._state_to_key(state)
    nk = self._state_to_key(next_state)
    cur_q = self.q_table[sk][action]
    if next_avail:
      max_nq = max(self.q_table[nk][m]
                   for m in next_avail)
    else:
      max_nq = 0  # game over
    # Q-learning formula
    new_q = (1-self.alpha)*cur_q + \
      self.alpha*(reward+self.gamma*max_nq)
    self.q_table[sk][action] = new_q
\end{lstlisting}

\subsection{RL Algorithm Categories}
\zhint{RL算法分类}

\textbf{1. Model-Free vs Model-Based}:
\begin{itemize}[leftmargin=8pt,itemsep=0pt,topsep=0pt]
\item \textbf{Model-Free}: learn policy directly (Q-Learning, DQN, Policy Gradient) \zhint{不建环境模型}
\item \textbf{Model-Based}: learn environment dynamics (AlphaZero, Dyna-Q) \zhint{先学环境再规划}
\end{itemize}

\textbf{2. Value vs Policy vs Actor-Critic}:
\begin{itemize}[leftmargin=8pt,itemsep=0pt,topsep=0pt]
\item \textbf{Value-Based}: learn V(s) or Q(s,a), derive policy (Q-Learning, DQN) \zhint{学价值函数}
\item \textbf{Policy-Based}: directly learn $\pi(a|s)$ (REINFORCE) \zhint{直接学策略}
\item \textbf{Actor-Critic}: hybrid of both (PPO) \zhint{结合两者}
\end{itemize}

\subsection{RL Evaluation}
\begin{enumerate}[leftmargin=10pt,itemsep=0pt,topsep=0pt]
\item Win rate change during training \zhint{训练过程中胜率变化}
\item Win rate vs random agent in testing \zhint{测试时对随机agent的胜率}
\item Number of different moves explored (Q-Table size) \zhint{Q表中探索的状态数}
\item Feedback from skilled opponents \zhint{专家评价}
\end{enumerate}

\end{multicols}
\end{CJK}
\end{document}
